\section{Question}
\begin{frame}{A question, a model, a comparison}
  \begin{columns}[T,onlytextwidth]
    \column{.57\textwidth}
    {\Large How does a decay rate\\change an oscillating response?}\par
    \vspace{6mm}
    \begin{enumerate}
      \item Define the variables and assumptions.
      \item Compare the response over time.
      \item Report a reproducible result.
    \end{enumerate}
    \column{.37\textwidth}
    \begin{block}{Scope}
      An analytic example for demonstrating document design. No experimental
      or simulation data are presented.
    \end{block}
    \begin{exampleblock}{One visual language}
      Clear headings, quiet backgrounds and a small set of consistent colors.
    \end{exampleblock}
  \end{columns}
\end{frame}

\section{Model}
\begin{frame}{Define the response before interpreting it}
  \begin{columns}[T,onlytextwidth]
    \column{.55\textwidth}
    \begin{equation*}
      \boxed{u(t;\gamma)=e^{-\gamma t}\cos t}
    \end{equation*}
    \vspace{4mm}
    \begin{itemize}
      \item $t\geq0$: dimensionless time.
      \item $\gamma>0$: decay rate.
      \item $u(0;\gamma)=1$: initial normalization.
    \end{itemize}
    \column{.39\textwidth}
    \begin{block}{Envelope}
      $A(t)=e^{-\gamma t}$ decreases monotonically.
    \end{block}
    \begin{alertblock}{Characteristic time}
      $\tau=1/\gamma$ sets the decay scale.
      The oscillation period remains $2\pi$.
    \end{alertblock}
  \end{columns}
\end{frame}

\section{Comparison}
\begin{frame}{Stronger damping reduces the envelope}
  \centering
  \pkuResponsePlot{.91\textwidth}{.62\textheight}\par
  {\footnotesize\color{pku-muted}Analytic curves. Markers and line patterns preserve series identity in grayscale.}
\end{frame}

\begin{frame}{Report quantities that can be checked}
  \begin{table}
  \centering
  \begin{tabularx}{.95\linewidth}{@{}Xrrr@{}}
    \toprule
    \pkuTableHead \textbf{Series} & $\gamma$ & $\tau$ & $A(6)$\\
    \midrule
    Slow decay & 0.12 & 8.333 & 0.487\\
    Moderate decay & 0.25 & 4.000 & 0.223\\
    \rowcolor{pku-tint} Fast decay & 0.45 & 2.222 & 0.067\\
    \bottomrule
  \end{tabularx}
  \end{table}
  \vspace{5mm}
  \begin{block}{Interpretation}
    Increasing $\gamma$ shortens the decay time. The table reports the
    exponential envelope, rather than the oscillating response itself.
  \end{block}
\end{frame}

\section{Reproducibility}
\begin{frame}[fragile]{Keep the implementation close to the result}
\begin{lstlisting}[language=Python]
import numpy as np

def response(t, gamma):
    return np.exp(-gamma * t) * np.cos(t)

time = np.linspace(0.0, 12.0, 121)
curves = {g: response(time, g) for g in (0.12, 0.25, 0.45)}
\end{lstlisting}
  \vspace{3mm}
  \begin{itemize}
    \item State the function, range and parameter values.
    \item Keep units and labels identical in every output.
    \item Replace illustrative examples with your own evidence.
  \end{itemize}
\end{frame}

\begin{frame}{Takeaways}
  \relax
  {\Large A small set of consistent decisions\\makes results easier to read.}\par
  \vspace{8mm}
  \begin{enumerate}
    \item Lead with the research question.
    \item Let figures and equations carry the argument.
    \item Use color to distinguish roles, not to fill space.
  \end{enumerate}
  \vfill
  {\color{pku-primary}\large Questions \& discussion}
\end{frame}
